\documentclass[12pt]{ca-capstone}
\geometry{margin=25mm}
\usepackage{setspace}
\setstretch{1.1}
\capstonerefs{crandall-pomerance,welschenbach,knuth-taocp2}
\begin{document}
\capstonetitle{Арифметика Монтгомери и быстрое возведение в степень}
\begin{center}\textbf{Теоретическая справка для преподавателя. Не выдавать студентам.}\end{center}

\section{Представление Монтгомери}
Пусть $n$ нечётно, $R=2^w>n$, $\gcd(R,n)=1$. Элементу $x$ соответствует $\tilde x=xR\bmod n$.
Выберем $n'=-n^{-1}\pmod R$. Для $0\le t<nR$ операция REDC вычисляет
\[
 m=t n'\bmod R,\qquad u=(t+mn)/R.
\]
Число $mn$ делится на $R$, $u\equiv tR^{-1}\pmod n$ и $0\le u<2n$, поэтому достаточно одного
условного вычитания. Следовательно,
$\operatorname{REDC}(\tilde a\tilde b)=abR\bmod n$. Перевод выполняется умножением на $R$,
обратный --$\operatorname{REDC}(\tilde x)$; монтгомериевская единица равна $R\bmod n$.
Для $n=1$ все результаты нулевые, но отдельная проверка в тестах обязательна.

\section{Бинарный и $q$-арный методы}
Бинарный метод слева-направо для каждого бита выполняет квадрат, а для единичного бита ещё и
умножение на основание. При $e=0$ цикл пуст, результат $1\bmod n$.

В $q$-арном методе $q=2^k$ заранее вычисляются $a^0,a^1,\ldots,a^{q-1}$. Для каждой $q$-арной
цифры $d$ выполняются $k$ квадратов и, если $d\ne0$, умножение на $a^d$. Нулевые старшие цифры
важны для числа квадратов; параметр $q$ должен быть степенью двойки не меньше двух.

\section{Скользящее окно}
Предвычисляются $a^2$ и нечётные степени $a^1,a^3,\ldots,a^{2^k-1}$. При сканировании битов
нули дают один квадрат. Начиная с единицы, выбирается окно длины не более $k$, оканчивающееся
единицей; его значение $d$ нечётно. Выполняются квадраты по длине окна и умножение на $a^d$.
Таблица индексируется самим нечётным показателем, а не его порядковым номером.

Для фиксированного основания таблица не пересчитывается. Ключ кэша обязан включать основание,
модуль, метод и параметр окна. Изменение любого из них при повторном использовании старой
таблицы является ошибкой корректности, а не только производительности.

\section{Число операций и benchmark}
При битовой длине $\lambda$ и числе единиц $\nu$ бинарный метод выполняет $\lambda$ квадратов
(включая возможный первый формальный квадрат) и $\nu$ умножений. Скользящее окно уменьшает число
умножений ценой $2^{k-1}-1$ предвычисленных нечётных степеней и дополнительной памяти. Поэтому
сравнение только времени одного запуска без учёта предвычисления некорректно.

Отдельно измеряются: построение контекста Монтгомери, перевод в домен и обратно, таблицы степеней,
первый запрос и повторные запросы. REDC-счётчик позволяет проверить, что все модульные умножения
действительно проходят через реализацию Монтгомери.

\clearpage
\section{Контрольный список}
Тесты: $e=0,1$, нулевое и некопримое основание, $n=1$, чётные и отрицательные модули, $a>n$,
окна разных размеров, нулевые внутри окна, повторные запросы и смена основания. Для REDC
проверяются тождество $t+mn\equiv0\pmod R$, диапазон результата и обратные переводы. Оракулом
служит трёхаргументный \texttt{pow}; трассировка сверяется с разложением показателя независимо.

\references
\end{document}
